% concatenated from main.tex
\documentclass[11pt]{amsart}
\usepackage[utf8]{inputenc}
\usepackage{csquotes}
\usepackage[foot]{amsaddr}


\usepackage[dvips]{graphics,epsfig}
\usepackage{amsmath,amsthm,amssymb}
\usepackage{array}
\usepackage{bbm}
\usepackage{xcolor}
\usepackage{verbatim}
\usepackage{url}
%\usepackage{authblk}
\pagestyle{plain}


\def\ve{\varepsilon}
\def\vf{\varphi}
\def\vn{\varnothing}

\def\AA{\mathcal{A}}
\def\AAA{\mathbb{A}}
\def\BB{\mathcal{B}}
\def\CC{\mathcal{C}}
\def\CCC{\mathbb{C}}
\def\DDD{\mathbb{D}}
\def\DD{\mathcal{D}}
\def\FF{\mathcal{F}}
\def\HH{\mathcal{H}}
\def\II{\mathcal{I}}
\def\III{\mathbb{I}}
\def\MM{\mathcal{M}}
\def\NN{\mathcal{N}}
\def\NNN{\mathbb{N}}
\def\PPP{\mathbb{P}}
\def\QQQ{\mathbb{Q}}
\def\RR{\mathcal{R}}
\def\RRR{\mathbb{R}}
\def\SS{\mathcal{S}}
\def\SSS{\mathbb{S}}
\def\UU{\mathcal{U}}
\def\VV{\mathcal{V}}
\def\WW{\mathcal{W}}
\def\ZZ{\mathcal{Z}}
\def\ZZZ{\mathbb{Z}}
\def\PP{\mathcal{P}}
\def\TTT{\mathbb{T}}
\def\EEE{\mathbb{E}}


%%%%  jb abbreviations
\def\bK{{\bf K}}
\def\bL{{\bf L}}
\def\bU{{\bf U}}
\def\aut{\mathop{aut}}
\def\betabar{{\bf\beta}}
\def\Uscr{\mathcal U}
\def\myreduct{\upharpoonright}
\def\abar{{\bf a}}
\def\bbar{\bf b}
\def\cbar{\bf c}
\def\lg{\mathop{\rm lg}}
\def\nmdeg{\mathop{\rm nmdeg}}
\def\icl{\mathop{\rm icl}}
% \sidebar{..} puts .. in comment form
\newcommand{\sidebar}[1]{\vskip10pt\noindent
 \hskip.70truein\vrule width2.0pt\hskip.5em
 \vbox{\hsize= 4truein\noindent\footnotesize\relax #1 }\vskip10pt\noindent}



 %%%%

\def\QQQa{\mathbb{Q}\langle \Bar{\alpha},\alpha_0 \rangle}

\def\cc{\mathfrak{c}}

\def\a{\alpha}
\def\b{\beta}
\def\g{\gamma}
\def\d{\delta}
\def\e{\varepsilon}
\def\z{\zeta}
\def\l{\lambda}
\def\k{\kappa}
\def\o{\theta}
\def\r{\rho}
\def\s{\sigma}
\def\t{\tau}
\def\w{\omega}
\def\D{\Delta}

\def\nmdeg{\operatorname{nmdeg}}
\newcommand{\Fmb}[1]{F_{Mb}\left( #1  \right)}
\newcommand{\indf}[1]{F_{ind}\left( #1  \right)}
\newcommand{\fcb}[1]{F_{cb}\left( #1  \right)}
\newcommand{\ficb}[1]{F_{Icb}\left( #1  \right)}





\DeclareMathOperator{\alg}{alg}
\DeclareMathOperator{\RU}{RU}
\DeclareMathOperator{\RM}{RM}
\DeclareMathOperator{\RD}{RD}
\DeclareMathOperator{\tp}{tp}
\DeclareMathOperator{\acl}{acl}
\DeclareMathOperator{\dcl}{dcl}
\DeclareMathOperator{\Cb}{Cb}
\DeclareMathOperator{\dtr}{\delta-trdeg}


\def\id{\operatorname{id}}
\def\Frac{\operatorname{Frac}}
\def\Const{\operatorname{Const}}
\def\spec{\operatorname{Spec}}
\def\span{\operatorname{span}}
\def\exc{\operatorname{Exc}}
\def\Div{\operatorname{Div}}
\def\cl{\operatorname{cl}}
\def\mer{\operatorname{mer}}
\def\trdeg{\operatorname{trdeg}}
\def\ord{\operatorname{ord}}




\newtheorem{theorem}{Theorem}[section]
\newtheorem{thm}[theorem]{Theorem}
\newtheorem{question}[theorem]{Question}
\newtheorem{lemma}[theorem]{Lemma}
\newtheorem{cor}[theorem]{Corollary}
\newtheorem{prop}[theorem]{Proposition}
\newtheorem{fact}[theorem]{Fact}
\newtheorem{claim}[theorem]{Claim}
\newtheorem{subclaim}[theorem]{Subclaim}
\newtheorem{conj}[theorem]{Conjecture}
\newtheorem{cond}[theorem]{Condition}
\newtheorem{observation}[theorem]{Observation}
\newtheorem{notation}[theorem]{Notation}

\theoremstyle{definition}
\newtheorem{definition}[theorem]{Definition}
\newtheorem{example}[theorem]{Example}
\newtheorem{xca}[theorem]{Exercise}
\newtheorem{ques}[theorem]{Question}
\newtheorem{note}[theorem]{Note}

\theoremstyle{remark}
\newtheorem{remark}[theorem]{Remark}
\newcommand{\m}{\mathbb }
\newcommand{\mc}{\mathcal }
\newcommand{\mf}{\mathfrak }
\newcommand{\mb}{F_{Mb} }
\newcommand{\G}{\mathrm{G} }
\newcommand{\U}{\mathrm{U} }



\def \ta {\tau_{\mathcal{D}/\Delta}}
\def \D {\Delta}
\def \DD {\mathcal D}

\newcommand{\gen}[1]{\left\langle#1\right\rangle}
%\renewcommand{\restriction}{\mathord{\upharpoonright}}

\newcommand{\pd}[2]{\frac{\partial #1}{\partial #2}}
\newcommand{\td}{\text{tr.deg.}}
\newcommand{\pp}{\partial }
\newcommand{\pdtwo}[2]{\frac{\partial^2 #1}{\partial #2^2}}
\newcommand{\od}[2]{\frac{d #1}{d #2}}
\def\Ind{\setbox0=\hbox{$x$}\kern\wd0\hbox to 0pt{\hss$\mid$\hss} \lower.9\ht0\hbox to 0pt{\hss$\smile$\hss}\kern\wd0}
\def\Notind{\setbox0=\hbox{$x$}\kern\wd0\hbox to 0pt{\mathchardef \nn=12854\hss$\nn$\kern1.4\wd0\hss}\hbox to 0pt{\hss$\mid$\hss}\lower.9\ht0 \hbox to 0pt{\hss$\smile$\hss}\kern\wd0}
\def\ind{\mathop{\mathpalette\Ind{}}}
\def\nind{\mathop{\mathpalette\Notind{}}}
\numberwithin{equation}{section}
\newcommand{\is}{^{p^ {-\infty}}}
\newcommand{\QQ}{\mathbb Q}
\newcommand{\fh}{\mathfrak h}
\newcommand{\SL}{\operatorname{SL}}


\title{No omega-stable type of limit Morley rank has nonminimality degree equal to the Morley-sequence invariant together with definable Morley sequences}
\author{John T. Baldwin, James Freitag, Scott Mutchnik}
\date{arxiv:2405.08211v3; CC BY 4.0}
\begin{document}
\maketitle
\setcounter{section}{7}
\setcounter{theorem}{27}
\noindent\textit{Editorial note.} The counted statement is the theorem labelled nmdegmbdefinablemorleyproperty at source lines 2887-2890, and the proof environment at source lines 2892-2897 is the proof of that same theorem by adjacency: it is the only proof environment attached to it, and the theorem is not restated anywhere else in the article, so the pairing needs no appeal to a named restatement.\par\medskip
\noindent\textit{Editorial note.} The retained context environment is the definition of the property named in the counted statement (definable Kim-Morley sequences) at source lines 2760-2762, which is also the definition introduced by this article. Because the theorem-like environments between it and the counted theorem are not retained, the section and theorem counters are set so that the counted theorem prints with the article's number Theorem 7.29; the retained definition therefore prints as Definition 7.28 instead of its article number.\par\medskip
\noindent\textit{Editorial note.} The article's own preamble is kept, including its amsart class, its theorem environments and its whole macro library. Only the times font package at source line 9 is left out, because it carries no TU-encoded font shapes under XeLaTeX and would silently substitute fonts; no macro is invented and no author text is changed.\par\medskip
\noindent\textit{Editorial note.} No retained line contains a citation command, a figure, an included graphics file or a reference to a label defined outside the retained spans, so the package declares no bibliography, no asset and no external reference, and no unresolved reference or citation can arise.\par\medskip
\noindent\textit{Editorial note.} The proof is a single dense paragraph. Every symbol, hypothesis and sentence of both the counted statement and its proof is delivered line for line exactly as printed, with no repair of notation and no added hypothesis.\par\medskip
\noindent\textit{Editorial scope.} One original result: the theorem labelled nmdegmbdefinablemorleyproperty (source lines 2887-2890) with its complete adjacent author proof (source lines 2892-2897), together with the author's own definition of definable Kim-Morley sequences (source lines 2760-2762). The appendix theorem at source lines 2941-2943 together with its proof at source lines 2945-2984, the finite-rank comparison of Corollary eqnmfin, and the invariants developed in Sections 3 to 6 are not counted here and are not used as premises.\par\medskip
\noindent\textit{Allowed standard prerequisites (editorial).}\begin{itemize}
\item Forking independence, Morley sequences and Morley rank in an omega-stable theory. For a complete type p over a set A, a Morley sequence in p over A is an A-indiscernible sequence of realizations of p every finite initial segment of which is independent over A; RM is Morley rank; over an algebraically closed base in an omega-stable theory every type is stationary and has a unique nonforking extension.. The first sentence of the proof passes to a base algebraically closed in T\textasciicircum{}eq so that p becomes stationary, and the last step identifies Morley-ness with independence of the n-th term from the preceding n-1 terms.
\item Kim independence and Kim's lemma for NSOP1 theories. Kim-independence is the independence notion of NSOP1 theories; Kim's lemma says that if a sequence is Morley over a base set, then a type over an extension of that base does not fork over the base.. Used at the end of the proof: were the constructed sequence a Morley sequence over A, Kim's lemma would make B independent from a\_0 over A, contradicting that the type of a\_0 over B forks over A.
\item Compactness and the strict gap below a limit ordinal. First-order compactness, including the extraction of a single formula from a type-definable family, and the elementary fact that below a limit ordinal lambda every ordinal lambda-prime below lambda has a strictly larger ordinal still below lambda.. Used twice: to extract the single formula phi from the set sigma of formulas that fork over A, and to choose the Morley rank lambda-prime-prime strictly between lambda-prime and the limit ordinal lambda.
\end{itemize}\medskip
\begin{definition}\label{definablemorleyproperty}
    Let $T$ be $\mathrm{NSOP}_{1}$ with nonforking existence, and let $C \subseteq \mathbb{M}$ be a (small) set, and let $p(x)$ be a complete type over $C$. Then $p(x)$ \textit{has definable $\ind^{K}$-Morley sequences} if the set of $\ind^{K}$-Morley sequences $\{b_i\}_{i < \omega}$ over $C$ with $b_{i} \models p(x)$ is definable over $C$ relative to the (necessarily type-definable) set of indiscernible sequences $\{b_i\}_{i < \omega}$ over $C$ with $b_{i} \models p(x)$. (That is, for $S_{1}$ the set of $\ind^{K}$-Morley sequences $\{b_i\}_{i < \omega}$ over $C$ with $b_{i} \models p(x)$, and $S_{2}$ the set of $C$-indiscernible sequences $\{b_i\}_{i < \omega}$ with $b_{i} \models p(x)$, there is a formula $\varphi(x_{1}, \ldots x_{n})$ with parameters in $C$, with $|x_{i}| = |b_{i}|=k$, so that $S_{1} = S_{2} \cap \{\{b_{i}\}_{i < \omega} \in (\mathbb{M}^{k})^{\omega} : \mathbb{M} \models \varphi(b_{1}, \ldots, b_{n}) \}$.)
\end{definition}

  \begin{theorem}\label{nmdegmbdefinablemorleyproperty}
Let $T$ be $\omega$-stable. Then there is no type $p \in S(A)$ such that $\mathrm{RM}(p)$ is a limit ordinal, $\mathrm{nmdeg}(p)= F_{Mb}(p)$, and $p$ has definable Morley sequences.

\end{theorem}

\begin{proof}
Because a set is independent over $A$ if and only if it is independent over $\mathrm{acl}^{eq}(A)$, and is an indiscernible sequence over $A$ if and only if it is an indiscernible sequence over $\mathrm{acl}^{eq}(A)$, the formula witnessing that $p$ has definable Morley sequences also witnesses, for any extension $p'$ of $p$ to $\mathrm{acl}^{eq}(A)$, that $p'$ has definable Morley sequences, and moreover, $\mathrm{nmdeg}(p') = \mathrm{nmdeg}(p) = F_{Mb}(p) = F_{Mb}(p')$ and $\mathrm{RM}(p) = \mathrm{RM}(p')$ is a limit ordinal. So we may assume $A$ is algebraically closed in $T^{eq}$, and therefore that $p$ is stationary. If $p$ has definable Morley sequences, it follows that the set of Morley sequences over $A$ consisting of realizations of $p$ is relatively definable in the set of \textit{nonconstant} $A$-indiscernible sequences consisting of realizations of $p$. Let $\mathrm{nmdeg}(p)= F_{Mb}(p) = n$, and for $q(y) = (p(x))^{\otimes{n}}$ the type of $n$ independent realizations of $p$ (which is well-defined by stationarity of $p$), let $\sigma$ be the set of formulas $\varphi(x, y)$ where $\varphi(x, a)$ forks over $A$ for any $a \models q(y)$. Since $\mathrm{nmdeg}(p)= F_{Mb}(p) = n$, a nonconstant $A$-indiscernible sequence $\{a_{i}\}_{i < \omega}$ is a Morley sequence if and only if $a_{n} \ind_{A} a_{0} \ldots a_{n-1}$, if and only if $\models \neg\varphi(a_{n}, a_{0} \ldots a_{n-1})$ for every $\varphi(x, y) \in \sigma$. By compactness, and the assumption that $p$ has definable Morley sequences over $a$, there is then some $\varphi(x, y)$ so that, for any $a_{0} \ldots a_{n-1} \models q(y)$, $\varphi(x, a_{0} \ldots a_{n})$ forks over $A$, and so that, for $\{a_{i}\}_{i < \omega}$ a nonconstant $A$-indiscernible sequence, $\{a_{i}\}_{i < \omega}$ is not an $A$-Morley sequence if and only if $\models \varphi(a_{n}, a_{0} \ldots a_{n-1})$. Fix some $b_{0} \ldots b_{n} \models q(y)$, and let $\mathrm{RM}(p) = \lambda$; we have assumed that $\lambda$ is a limit ordinal. Then $\varphi(x, b_{0} \ldots b_{n})$ forks over $A$, so $\lambda' =: \mathrm{RM}(\varphi(x, b_{0} \ldots b_{n}) \cup p(x)) < \lambda$, because otherwise there would be a completion $\tilde{p}(x)\in S(Ab)$ of $\varphi(x, b_{0} \ldots b_{n}) \cup p(x)$ so that $\mathrm{RM}(\tilde{p(x)})=\mathrm{RM}(\varphi(x, b_{0} \ldots b_{n}) \cup p(x)) = \mathrm{RM}(p(x))$, implying that $\tilde{p}(x)$ does not fork over $A$ despite containg $\varphi(x, b_{0} \ldots b_{n})$, a contradiction. Since $\lambda$ is a limit ordinal, there is an ordinal $\lambda''$ so that $\lambda' < \lambda'' < \lambda$. There is then some $q(x) \in S(B)$ extending $p(x)$ so that $\mathrm{RM}(q(x)) = \lambda''$. Let $\{a_{i}\}_{i < \omega}$ be a Morley sequence in $q(x)$ over $B$. Then $\{a_{i}\}_{i < \omega}$ is nonconstant, but is not a Morley sequence over $A$, because if it were, by Kim's lemma $B \ind_{A} a_{0}$, contradicting that $a_{0} \models q(x)$, which forks over $A$. However, $\mathrm{RM}(a_{n}/a_{0} \ldots a_{n-1}) \geq \mathrm{RM}(a_{n}/Ba_{0} \ldots a_{n-1}) = \mathrm{RM}(a_{n}/B) = \mathrm{RM}(q(x))= \lambda''$, the first equality because $\{a_{i}\}_{i < \omega}$ is a Morley sequence over $B$. However, because $\{a_{i}\}_{i < \omega}$ is nonconstant and not a Morley sequence over $A$, $\varphi(x, a_{0} \dots a_{n-1}) \in \mathrm{tp}(a_{n}/a_{0} \ldots a_{n-1})$, so $\mathrm{RM}(\varphi(x, a_{0} \ldots a_{n}) \cup p(x)) \geq \mathrm{RM}(a_{n}/a_{0} \ldots a_{n-1}) \geq \lambda''$. But $a_{0} \ldots a_{n-1} \equiv_{A} b_{0} \ldots b_{n-1}$, so $\lambda' = \mathrm{RM}(\varphi(x, b_{0} \ldots b_{n}) \cup p(x)) = \mathrm{RM}(\varphi(x, a_{0} \ldots a_{n}) \cup p(x)) \geq \lambda'' > \lambda'$, a contradiction.



\end{proof}

\end{document}
